% Tipo de documento. En este caso es un art�culo, para folios A4, tama�o de la fuente 11pt y con p�gina separada para el t�tulo
\documentclass[a4paper,12pt,titlepage]{article}

% Carga de paquetes necesarios. OrdenesArticle es un paquete personalizado
\usepackage[spanish]{babel} 
\RequirePackage[T1]{fontenc}
\RequirePackage[ansinew]{inputenx} 
\usepackage[spanish,cap,cont,title,fancy]{OrdenesArticle}
\usepackage{array}
\usepackage{graphicx}
\usepackage{hyperref}
\usepackage{pifont}
\usepackage{listings}
\usepackage[usenames,dvipsnames]{color}
\usepackage{colortbl}
\usepackage{makeidx}
\hypersetup{bookmarksopen,bookmarksopenlevel=3,linktocpage,colorlinks,urlcolor=blue,citecolor=blue,
						linkcolor=blue,filecolor=blue,pdfnewwindow,
						pdftitle={Listado de preguntas Funcionalidad 1}, 
						pdfauthor={Juan Andrada Romero, Juan Jos� Antequera Flores},
						pdfsubject={Modelos Avanzados de Base de Datos}}


% Macro para definir una lista personalizada 
\newenvironment{milista}%
{\begin{list}{\textbullet}%
{\settowidth{\labelwidth}{\textbullet} \setlength{\leftmargin}{\dimexpr\labelsep+\labelwidth+5pt}
\setlength{\itemsep}{\dimexpr 0.5ex plus 0.25ex minus 0.25ex}
\setlength{\parsep}{\itemsep}
\setlength{\partopsep}{\itemsep}
\addtolength{\topsep}{-7.5pt}
}}%
{\end{list}}


\begin{document}

% En las p�ginas de portada e �ndices, no hay encabezado ni pie de p�gina
\pagestyle{empty} 

% Se incluye la portada
\thispagestyle{empty}
\begin{center}
  {\LARGE UNIVERSIDAD DE CASTILLA-LA MANCHA} \\
  \bigskip
  {\Large ESCUELA SUPERIOR DE INFORM�TICA} \\
  \vspace{28mm}
  \includegraphics[scale=0.45, keepaspectratio]{./imagenes/esi_bw.png} \\
  \vspace{30mm}
  {\LARGE \textbf{Modelos Avanzados de Base de Datos}} \\
  \vspace{14mm}
  {\Large \textsf{\textsc{Listado de preguntas del trabajo Distribuci�n 1:}}} \\
  \vspace{2mm}
  {\Large \textsf{\textsc{Bases de Datos Distribuidas}}}\\
  \vspace{15mm}
  {\large Juan Andrada Romero} \\
  {\large Juan Jos� Antequera Flores} \\
  {\large Jose Domingo L�pez L�pez} \\
  \vspace{9mm}
  {\large \today}
\end{center}
\clearpage

\parskip = 10pt
% Comienza el contenido del documento. Se utilizan n�meros ar�bigos y el encabezado y pie de p�gina personalizado
\pagenumbering{arabic}
\pagestyle{fancy}

\paragraph{Pregunta:} En el trabajo se comentan los lenguajes de definici�n de objetos y de consulta de objetos que define el est�ndar ODMG, pero no se hace referencia a ning�n lenguaje para poder manipular objetos, es decir, insertarlos en la base de datos, modificarlos o eliminarlos. �Dicho est�ndar define un lenguaje de manipulaci�n de datos espec�fico? �C�mo se realizar�an las manipulaciones sobre objetos en una base de datos orientada a objetos?

\paragraph{Respuesta:} El est�ndar ODMG-93 sugiere que el lenguaje de manipulaci�n de objetos (OML) sea la extensi�n de un lenguaje de programaci�n, de forma que se pueden realizar entre otras las siguientes operaciones sobre la base de datos: creaci�n, borrado, y modificaci�n de un objeto. \cite{bdoo}.

Por tanto, los objetos se manipulan dentro del lenguaje de programaci�n orientado a objetos, utilizando el concepto de transacci�n y de clases persistentes. De este modo, como indica el ODMG API (\cite{odmgapi}), para realizar las operaciones de inserci�n, modificaci�n o eliminaci�n de un objeto dentro del lenguaje de programaci�n, se proceder�a de la siguiente forma:
\begin{enumerate}
	\item Crear una transacci�n.
	\item Obtener la base de datos en uso.
	\item Ejecutar la operaci�n de inserci�n, modificaci�n o eliminaci�n.
\end{enumerate}

Un ejemplo de inserci�n ser�a el siguiente:

\begin{verbatim}
 public static void storeNewProduct(Product product)
	    Transaction tx = odmg.newTransaction();
	    tx.begin();
	    // get current used Database instance
	    Database db = odmg.getDatabase(null);
	    // make persistent new object
	    db.makePersistent(product);
	    tx.commit();
\end{verbatim}




\newpage
% Bibliograf\'ia 
\bibliographystyle{plain}
\pagestyle{plain} 
\bibliography{Bibliografia}


\end{document}

